\subsection{Packaging}\label{sub:Packaging} % (fold)

L'encodeur pré-entraîné est distribué sous la forme d'un paquet Python installable,
nommé \texttt{elf64\_context\_hash}, conçu selon le principe du \textit{zero-friction
deployment} : l'utilisateur n'a à connaître ni l'architecture du modèle, ni les chemins
de fichiers, ni le protocole d'extraction VEX — l'ensemble du pipeline est encapsulé
derrière une interface uniforme.

\subsubsection{Structure du paquet et point d'entrée CLI}

Le paquet est déclaré via un fichier \texttt{pyproject.toml} standard et expose un
unique point d'entrée en ligne de commande :

\mintinline{sh}{elf64ctx [-E <binaires>] [-C <f1> <f2>] [-P] [-F <fn1> <fn2>] [--plot] [-D]}

Les deux modes principaux sont l'encodage (\texttt{-E}) et la comparaison
(\texttt{-C}). En mode encodage, le pipeline complet est exécuté sur chaque fichier ELF
fourni : extraction du Bag-of-Paths via \texttt{angr}, canonicalisation VEX par la
projection $\Phi$, encodage par le \texttt{Conv1DEncoder}, normalisation $\ell_2$, et
sérialisation des embeddings par fonction au format JSON (pickle + encodage Base64). En
mode comparaison, les embeddings préalablement sérialisés de deux binaires sont chargés
et une matrice de distances euclidiennes dans $\mathbb{R}^{256}$ est calculée pour toutes
les paires de fonctions, avec option de visualisation en heatmap (\texttt{--plot},
colormap \texttt{plasma\_r}).

Le flag \texttt{-P} active un mode prédicteur : un token \texttt{<MASK>} est annexé à
chaque séquence avant encodage, et c'est la sortie du prédicteur $g_\varphi$ à la
position finale — plutôt que la sortie brute de $f_\theta$ — qui est retenue comme
embedding. Ce mode permet d'obtenir une représentation prospective, conditionnée sur la
continuation latente attendue du chemin d'exécution.

\subsubsection{Téléchargement automatique des artefacts depuis HuggingFace Hub}

Le vocabulaire canonique (\texttt{vocab.json}, $|V| = 371$ tokens) et le checkpoint
d'entraînement (\texttt{latest.pt}, contenant les états de $f_\theta$ et $g_\varphi$)
sont hébergés sur le dépôt HuggingFace \href{https://huggingface.co/rubenanko/binary-jepa/tree/main}{\texttt{rubenanko/binary-jepa}}. Leur
téléchargement est déclenché automatiquement lors du premier import du module, via
\texttt{hf\_hub\_download}, dans le répertoire \texttt{\textasciitilde/.elf64-context-hash/} :

\begin{minted}{Python}
if not DEFAULT_VOCAB_PATH.is_file():
    hf_hub_download(repo_id=DEFAULT_REPO_ID,
                    filename="vocab.json",
                    local_dir=DEFAULT_DATA_PATH)
\end{minted}

Ce mécanisme garantit la reproductibilité : toute installation du paquet, quelle que
soit la machine, converge vers le même checkpoint de référence sans intervention
manuelle. Il découple par ailleurs le cycle de vie du modèle — mis à jour sur le Hub à
chaque ré-entraînement — de celui du paquet, versionné indépendamment.

\subsubsection{Isolation des dépendances lourdes}

Le module charge \texttt{torch}, \texttt{angr} et \texttt{matplotlib} à la demande,
dans les fonctions qui en ont besoin, plutôt qu'au niveau du module. Cela maintient un
temps d'import minimal pour les usages programmatiques qui n'auraient besoin que du
chargement des embeddings pré-calculés (\texttt{load\_embeddings}), sans déclencher
l'initialisation coûteuse d'angr ou le chargement des poids PyTorch. La séparation
entre \texttt{loaders.py} (chargement des artefacts), \texttt{elf\_processing.py}
(extraction statique) et \texttt{model/} (architectures) reflète cette logique de
dépendance paresseuse.
% subsection Packaging (end)
